%% The researcher's edit of the practitioner section (HumanEdited/practitioner.tex), restructured by Claude
%% (2026-10-23): one subsection per point, blocks moved as the red notes asked, references filled, attributions
%% restored. Sentences Claude added are blue; replies and questions are blue and in [Claude: ...].

% ======== q-practitioner ========

\subsection{A practitioner's reading}\label{sec:qual:practitioner}


One participant who signed up had a history of experience in the British Army, having completed Officer Training as well as operational tours. His experience with real-life strategic and tactical command, dispatch, and working with human subordinates in high-stress scenarios made him an interesting case study. He completed the experiment exactly the same as the others, and then gave a more in-depth interview. 

\textcolor{blue}{[Claude: Restructured into ten subsections, one per point, in place of the run-in paragraphs. The labels the rest of the paper cites are kept on the subsection they now belong to. Blocks your notes said belong elsewhere (the pace, the parallel-work and skills quote, and the old ``read against the results'' list) are held at the end under ``Held for other sections''; the red-team agent has moved to the queue-ordering section of qual.tex. Attributions now say only the source, with no timecodes, and nowhere says ``includes the researcher's words''.]}

\subsubsection{How typical he was}\label{sec:qual:practitioner:own-play}

Figure~\ref{fig:qual:practitioner:own-play} compares our practitioner with the rest of the cohort to determine if he was markedly different to the civilian population. He was a good performer (0.7 SD above the game average, the 80th percentile) with a slightly higher work rate than most: quicker at making strategic decisions (9.6\,s per decision against a median of 11.4\,s) and more often deciding before the strategy cards had loaded (32\% of decisions against a median of 10\%). He was otherwise typical, sitting in the middle half of the cohort.

\begin{figure}[H]
\centering
\includegraphics[width=\linewidth,height=0.50\textheight,keepaspectratio]{figs/results/q-practitioner__own-play__play.png}
\caption{\textbf{Typical play.} One row per measure, each on its own scale: grey dots are the 33 participants, the shaded band the middle half, the tick the median and the large dot our practitioner.}
\label{fig:qual:practitioner:own-play}
\end{figure}
\textcolor{blue}{[Claude: The ``Before the study'' rows are gone (new figure: \texttt{q-practitioner\_\_own-play\_\_play.png}); the caption no longer says ``unusual background''.]}

\subsubsection{Decisive, and playing for the whole board}\label{sec:qual:practitioner:pt-whole-board}\label{sec:qual:practitioner:pt-quicker-by-hand}\label{sec:qual:practitioner:pt-resource-rich}

What was immediately apparent was his decisiveness. He worked quickly and handled the stress well, seldom being sucked into a single task and often returning to the broader picture, stating the importance of how ``success is generated across all of the tasks rather than just an individual task''. Like others, he relied on the agent more for complex problems, and manually pre-empted on smaller tasks often: 

\begin{quote}\small
``Do you know what? Probably quicker just to hit three on there. But we'll go three on that one there. [\ldots{}] I probably should have done that one manually. It would have been quicker.''\par\hfill--- P36 (think-aloud, session 1)
\end{quote}

\subsubsection{Simple by hand, complex to the assistant: the branching space of failure}\label{sec:qual:practitioner:pt-branching-paths}\label{sec:qual:practitioner:pt-resource-cascade}

Section~\ref{sec:res:size-decides:speed} gives a few perspectives as to why automation is preferred for the more complex tasks. Our practitioner agreed with the common logic of speed:

\begin{quote}\small
``sometimes it's going to be a little bit quicker to do these myself only because I know the answer based on what I'm seeing visually in front of me so i don't need to do that much [...] simple task you probably want to do yourself if you're going to get the solution quicker and the more complex tasks you're going to want somebody to give you that solution in front of you so you can give it a quick once over and go that looks about right''{\color{blue}\par\hfill--- P36 (think-aloud, session 1; continues below)}
\end{quote}

For the tactical agent specifically, he framed this in a different way to the rest. He was automatically planning to handle failure; his allocations considered setting up his next actions if that task were to fail, and when a task was sufficiently complicated he gave up attempting to anticipate this with a direct plan. Essentially, he saw larger tasks as complex enough that it was more important to get started and fix issues later:

\begin{quote}\small
``I would have to go and do a much more thorough scan of the [agent's] solution to say 'have you put redundancy into that plan or do i care'. I've put redundancy into that plan so i don't have to come back to it. [...] If one of my drones fails my assumption is that my other drones are picking up the task whereas the more complex ones just [I am] literally looking 'have they allocated enough to [get] that task started going because if it fails; think about hypertextual branching paths\footnote{{\color{blue}In hypertext fiction each choice branches into a different continuation of the story; by analogy, each possible failure in a complex mission branches into a different plan that would then be needed.}}: so as soon as I get one failure, that's going to change the entire story. [...] so it's going to need more human intervention to go into it so there's no point in me trying to think about all of that ahead of time. I may as well let them start on that task and then intervene when i need to''\par\hfill--- P36 (think-aloud, session 1)
\end{quote}

{\color{blue}He saw the same coupling between the tiers when asked about automating one level and not the other: whatever is allocated at the top sets the risk carried below.}

\textcolor{red}{Check this quote again against the audio.}

\begin{quote}\small
{\footnotesize\textit{Researcher:} one thing that you did clearly is you very much immediately were happy to sort of automate the low level of detail the tactical level and spend more of your attention on the high level which makes a lot of sense to me and that's how probably about half my participants have behaved half the participants have said well I want to automate the high level and then I can fix it}\par\smallskip
``I would argue that if you've not got the resources available to you at the low level then or you've over allocated resources at the high level then you're running out of resources cascade down so you're it's a balance it's never perfectly weighted one way or the other so if you're let's say you split the task this dude here is doing this task and all of a sudden has got no resource left look at the task I need more resource I told you need more resource when I told you this is my plan and this is I've got no redundancy here you're carrying risk whereas if you've done both it's your risk to carry on the flip side of it if you've over allocated a resource that's just wasted resource now I can't complete you might want more resource but I've got five tasks that I need to do so I'm balancing these five do what you can with what you've got let me know where it all falls down because it doesn't matter if yours falls down because yours is a moderate task I've got a critical one over here that I'm allocating more to''\par\hfill--- P36 (debrief)
\end{quote}

\subsubsection{Named rules and priorities}\label{sec:qual:practitioner:pt-reserve}\label{sec:qual:practitioner:pt-critical-first}

Whether consciously or not, our practitioner identified rules and adhered to these, such as ``never deploy without a reserve''. A named rule of this kind may be one way the consistency of individual styles (Section~\ref{sec:res:personal-style}) could be built through training, which we will discuss later in Section~\ref{sec:discussion}. He also returned to the use-case when considering task prioritisation:

\begin{quote}\small
``I've run out of blue so that's gonna hurt my score but at the same time I've got the critical tasks done. [...] If this is a real life scenario I'd rather be focused on the critical tasks than getting those done to a higher value rather than moderate tasks of getting them done to a low''\par\hfill--- P36 (think-aloud, session 1)
\end{quote}

\textcolor{blue}{[Claude: Training has no subsection of its own yet, so this points to the Discussion as a whole.]}

\subsubsection{Trusting the Strategic Assistant more}\label{sec:qual:practitioner:pt-strategic-trust}

Rarely amongst the cohort (Section~\ref{sec:res:tier-trust:gap}), he trusted the Strategic Assistant more than the Tactical. His reasoning for this is mostly about the consequences of a bad plan from the tactical being slightly greater: a mission might not be covered, whereas for the strategic it was just less optimal. It is worth noting (not shown here for brevity) that he does slightly caveat this position by then pointing out that a poor tactical agent can cause knock on effects later down the line.

\begin{quote}\small
``the strategic one I was less concerned about the solutions it presented so I hesitated less on the strategic than the tactical but only because all I was looking at there was: 'have I got a reserve allocated' whereas [for] a tactical one it's 'have you covered all the tasks' [\ldots{}] I wouldn't trust any of them without checking but I trust the strategic slightly more because the consequences are slightly less''\par\hfill--- P36 (after session 1)
\end{quote}

\subsubsection{Assistants as staff}\label{sec:qual:practitioner:pt-staff-process}

We asked the practitioner to consider his experience of command structure on this problem; military command builds hierarchies of concern with each person having a clear set of roles and responsibilities, and the relationship between them is crucial to the overall operation. We are proposing the incorporation of AI systems in structures of this type for search and rescue, so much can be learned from the comparison. When asked about the parallels between commanding two subordinates and commanding two agents:

\textcolor{blue}{[Claude: On the long quote: split as you have it, with the three parts marked as one continuous answer. The first part says it begins a continuous answer, the middle one says (continued), and the last carries the full attribution. Each commentary paragraph sits between them, so the three points read as one train of thought.]}

\begin{quote}\small
``well it's kind of this how it works in military planning anyway so in anything above battlegroup level [our 'strategic'] there would be a staff there that would work up plans, socialise those plans amongst themselves, and back-brief me as a commander''{\color{blue}\par\hfill--- P36 (debrief; the first of three consecutive extracts from one answer)}
\end{quote}

This justifies strongly the nature of agents working at the strategic and operational planning level as independent planners whose work is checked by a human overseer, as our study does at both levels. 

\begin{quote}\small
``I would then go 'okay yeah I want you a bit of that plan and a bit of that plan give me this new plan'. [...] So it is always a collective input into it [...] because I can't as an individual work through: the ground, the enemy, the enemy's plan, the likelihood of my plan succeeding, the logistics behind it. It would take me forever which is why your staff do it, so at that level you're already in a position where somebody's just coming back to you with the answers''{\color{blue}\par\hfill--- P36 (debrief, continued)}
\end{quote}

It also introduces the need for a human to be able to merge parts of multiple different plans: a possible future direction is a set of agents building parts of the plan from different perspectives, and a human-guided agent which can then merge these into a resultant plan.

\begin{quote}\small
``[and] as long as you have trust in how they've come up with those answers then if it looks right [gestures to 'go ahead']: because you've gone through the process yourself [...] you'd have worked up in those levels so you kind of understand how those answers are generated so there's a level of trust in that process''{\color{blue}\par\hfill--- P36 (debrief, end of the same answer)}
\end{quote}

Finally, it reinforces the importance of trust. Many participants spoke about trusting the thought process of the agent (Section~\ref{sec:qual:trust-unchecked}) by understanding it. The above quote invokes the idea of ``you'd have worked up in those levels so you kind of understand how those answers are generated''; this is something that is true of the human-human dyad, but not the human-agent. This perhaps is the most important difference when comparing the two systems.

\subsubsection{Socialising the plan, and how much context each level needs}\label{sec:qual:practitioner:pt-trust-but-verify}\label{sec:qual:practitioner:pt-silos}\label{sec:qual:practitioner:pt-one-up-two-down}

He described ``socialising the plan'' through taking it to each level of the hierarchy and giving every level an opportunity to find and highlight flaws:

\begin{quote}\small
``I write a plan I will go and take it to the lowest tactical level (which is something like this where you're moving different units around). You'd probably speak to the platoon commander, you'd speak to the platoon sergeant: 'I'm thinking of doing this'. You'd bring your section commanders in, you'd go through yours and that's the process: 'guys this is what I'm thinking and this is what your likely mission is, go away and work out the plan'. [...] As humans we should be used to socialising our decision making process to an extent. How you do that is entirely up to the commander but at the same time I'm still going through a process where somebody is checking my thinking and has the option to challenge me. So with AI, it should always be open to challenge [...] but also trust in the fact it's done the maths behind it and has allocated enough platforms to do the task you're after''\par\hfill--- P36 (think-aloud, session 2)
\end{quote}

\textcolor{blue}{[Claude: Section~\ref{sec:qual:trust-unchecked} is the nearest kept section, about trusting what cannot be checked. The material on operators wanting to know how the assistants decide was in two places you have since dropped: the pruned ``Trust in the process'' finding, and ``What operators asked for'', where 12 of 33 asked or wondered how the assistants decided. Say if you want a sentence back for it.]}

{\color{blue}The same socialising reaches down the hierarchy, and with it the question of how much each level needs to know. Narrow context is normal in staff planning, and fine as long as each part has been through a sound analytical process. This sits directly against the operators who wanted one plan held across both tiers (Section~\ref{sec:qual:hierarchy:one-plan}).}


\begin{quote}\small
{\footnotesize\textit{Researcher:} both of them are myopic in terms of what they see so they see only the current thing so the strategic agent all it sees is the list of drones and it has a simple rule that says one drone for each dot plus one of each or something like that the fact that they see only that context how would you say that compares to the scenario you're talking about with humans planning across multiple layers this might be a six or seven layer command structure}\par\smallskip
``You kind of work in your own silo so if you're the g2'er [intelligence branch of general's staff planning team] you just focus on the enemy and you're delivering what the enemy is likely to do and you have a team working to you to do that so as a g2 officer you'll have like a one person will work on the enemy capability, one person will look at on the ground, one person will look on the sort of the human terrain and they'll work through their own analytical processes to come up with an answer then gets aggregated up and you go into it so going into that planning brief you as a g2'er are will only see your part of the plan and then the g3 [operations branch] guys will see their part of the plan and then that's kind of the first time you come together but generally speaking you're you're not getting all of the context okay so but you don't that you need that because as a person you need to be able to work through the context solution part of it yeah and this type of scenario here you wouldn't need that because of the simplicity of the task it's i've got to go and do these things and for those things i need these allocations so it's hard to sort of translate across the two as long as in each of them it's gone through the analytical process and spat out an answer''\par\hfill--- P36 (debrief)
\end{quote}

\noindent\textit{Bears on:}
\begin{itemize}[leftmargin=1.2em,itemsep=1pt]\small
\item \textit{Qualifies} Section~\ref{sec:qual:hierarchy:one-plan}, \emph{One plan across both decisions}: Several operators wanted one plan held in mind across the tiers. P36 sees narrow context as normal (``you as a g2'er are will only see your part of the plan'') and fine ``as long as in each of them it's gone through the analytical process'', and adds that this game's tasks are simple enough not to need more: a direct tension with those operators.
\item \textit{Goes beyond} Section~\ref{sec:qual:hierarchy:one-plan}, \emph{One plan across both decisions}: Operators wondered how the two assistants fit together. P36's model is staff cells that each see their part and first ``come together'' at the planning brief.
\end{itemize}

{\color{blue}Each level, in turn, should be given the part of the level above that its task needs, and no more; an agent that reaches past its remit gives worse answers.}


\begin{quote}\small
{\footnotesize\textit{Researcher:} there's a philosophy of understanding what's going on one level above you and one level below you to have that or one or two levels above you or something for the context right so you you get told to go and i don't know um recce that road you need to know why you're wrecking that road so there's something above right and the idea that actually you need to have each of your agents in this in this scenario understanding the context a little above them i understand what you're saying that you don't want them to be thinking about the super wide context but is there an advantage to that linking up of understanding what's a little way up above you}\par\smallskip
``so it depends yeah so you could have your agent that's checking the colors might be the one that understands two up [the intent of those above you, and one about that] so you understand two up in the same way as you plan two down issue orders one down so you think about i'm going to be sure to report to my section commanders but my i'm thinking about what each of those individual soldiers doing with that scenario and that might be your agent that corrals all your thing but in a when we go through that staff process the person who's just doing a scenario they will understand the one up they will understand or part of the one up that's given to them as their task is there's an area of responsibility there's an area of interest there's a thing that the enemy is doing within this space i just need to kind of understand that that so it might be you might understand that your mission that's come to you within your organisation so that's technically your organisation's mission but you might consider that your one up so it's depending on on as long as the person doing the task is given the understanding they're required to do that task then it doesn't matter how it docks into everything else that's not their job that's the person above them to do that so for your agents yes it might be useful but no it might not be useful it depends and particularly it's not useful if they're giving shit answers like if they're coming out and going no i've allocated five over here because there's a thing over there it's like cool that's not your job your job is to do this thing over here and then there's other agents to do this thing over here and they don't need them to do both because that's where this agent is going to do both for them and that's quite often even humans doing it humans start thinking or g2 again they'll start thinking blue plan but your job is not to do the blue plan your job is to do the red plan stop thinking about or the blue is probably going to do this so i'm going to do this no no but i need to know that sir because i need to understand the context of what the red is going to do so you need to have an idea of what blue capabilities have but you don't have control over the blue plan and you're not there to anticipate the blue plan you're there to put it within this context just stick to your lane and do that output because someone else will do that for you''\par\hfill--- P36 (debrief)
\end{quote}

\noindent\textit{Bears on:}
\begin{itemize}[leftmargin=1.2em,itemsep=1pt]\small
\item \textit{Goes beyond} Section~\ref{sec:qual:hierarchy:one-plan}, \emph{One plan across both decisions}: How should the tiers fit? From doctrine: ``you understand two up in the same way as you plan two down'', each task comes with the part of the level above that it needs, and an agent reaching outside its lane is told ``that's not your job''.
\item \textit{Qualifies} Section~\ref{sec:qual:hierarchy:one-plan}, \emph{One plan across both decisions}: The operators who wanted one shared plan are answered by a different design: each level gets ``the understanding they're required to do that task'', and ``it doesn't matter how it docks into everything else''.
\end{itemize}

\subsubsection{Losing the plan you did not build}\label{sec:qual:practitioner:pt-play-catch-up}

One concern we identified is that delegation to AI removes the mental model a human operator may have of the plan. We capture this to some extent in Section~\ref{sec:res:size-decides:tactical}, and our expert agrees:

\begin{quote}\small
``if I did [all the allocations] manually I would know all of that because I've built the plan I've got it in my head. If I just said 'go nuts', when something goes wrong I've got to play catch up.''\par\hfill--- P36 (think-aloud, session 2)
\end{quote}

\subsubsection{Failure rates need context}\label{sec:qual:practitioner:pt-failure-rate}

\textcolor{red}{Also really good. I think this also links to training. I'm starting to think we put training solely in discussion and then we've found many little links to it.}

A common request from our main cohort (who were mostly AI experts and academics) was for data about the expected failure rates. Many asked directly or suggested that the interface should show the prior probability of the agent failing. Our expert tentatively agrees with this, but highlights a very important issue with people's perception of such a number:
\begin{quote}\small
``That always messes with people's heads doesn't it? Because we look at the 40\% failure chance [and think] no chance [you'd use/trust it] [...] but if you put a person into that scenario and it was all people doing [every role], your failure chance would be this as well [...] in the search and rescue environment you're always gambling with someone's life, like the failure rate of doing nothing is 100\% so you've got to do something. You're just trying to chip away at it and prioritise accordingly so having that information presented [is] great but it's then down to that person that you're training''\par\hfill--- P36 (debrief)
\end{quote}

\subsubsection{Real stakes and shared risk}\label{sec:qual:practitioner:pt-shared-risk}

When asked about the decisiveness of his action, the practitioner attributed this to experience but also to knowing it was not a real scenario, but brought up the level of comfort one builds up in high-stress scenarios when lives are at risk:

\begin{quote}\small
``I know it's a simulation. There isn't someone's life on the line at the end of it and I've done enough planning activities to know and I'm very visual so looking at a glance: 'yes it's allocated what I needed to allocate it cool let's start running with it' whereas I can imagine if you've not had that experience that I've generated through non-use of drone swarming technology - if you've not had that, if you're a young platoon commander who's literally just come out of Sandhurst\footnote{UK Army Officer Training School} [...] either they're gonna be foolhardy and just f***** go for it and people are gonna die, or they'll be almost paralysed because it's like 'oh s***, no if I do that someone is gonna die here's 40\% failure rate on this task''{\color{blue}\par\hfill--- P36 (debrief; continues below)}
\end{quote}

He then reinforces again the importance of socialisation, this time not in spotting faults, but in making well-informed decisions with people's lives.
\begin{quote}\small
``that's always been true regardless of whether or not technology was in there but before there was a level of interpersonal contact: 'I'm out there with my people. I might die alongside them so when I'm putting them under that risk [it's] because I'm taking the same risk myself, whereas something like this where you're not necessarily taking the risk yourself and [...] you've put humans in the way then that's going to then come back and impact on your decision-making. So I imagine you're going to see me less decisive if someone's gonna die at the end''\par\hfill--- P36 (debrief)
\end{quote}

\subsubsection*{Held for other sections}\label{sec:qual:practitioner:held}

\textcolor{blue}{[Claude: Kept here so nothing is lost, each where your notes said it belongs. Delete each once it has moved.]}

\paragraph{For the Limitations section: physics sets the pace.}\phantomsection\label{sec:qual:practitioner:pt-physics-pace}\label{sec:qual:practitioner:pt-tempo-echelon}

The pace of the game is not realistic: in a real operation physics bounds the pace, so decisions get faster but the operation does not, and the operator's load moves to waiting and contingency planning.

\begin{quote}\small
{\footnotesize\textit{Researcher:} so I would ask in your professional opinion obviously this is clearly we're trying to get at something that we think is of interest in this kind of setting and stuff in reality I'd imagine the strategic and tactical thing would probably be done by two people but how relevant would you say this kind of set up is}\par\smallskip
``so in terms of delivery absolutely in terms of the pace probably not so much [...]if you think so the biggest thing for us that is always hard to convey to decision makers is: physics remains the bounding factor I've got to get from this point here to that point there and if the adversary can see me from 20k away then I'm doing that activity from 20k and if you're going at 3k an hour then that's going to take you a long time to get from here to over there so there's a level of that acceleration in terms of decision however on the flip side of it when the shit hits the fan you are probably going to be very frustrated; As I move this thing to here I've now got to wait for this thing to get into place to go and do this thing here hurry up and get there but at the same time as a commander you're sat there thinking okay what am I going to do if this thing fails what am I then going to do what am I going to do so your mental load is still loaded it's just loaded in a slightly different way so you're always trying to balance cognitive capacity against the time pressure effectively''\par\hfill--- P36 (debrief)
\end{quote}

\noindent\textit{Bears on:}
\begin{itemize}[leftmargin=1.2em,itemsep=1pt]\small
\item \textit{Qualifies} Section~\ref{sec:res:workload:raw-vs-adjusted}, \emph{Queue depth against reliance --- raw, and adjusted}: A busier queue did not push anyone onto the assistants. P36 puts the real pressure elsewhere: ``physics remains the bounding factor'', and the load becomes waiting and asking ``what am I going to do if this thing fails''. ``Your mental load is still loaded it's just loaded in a slightly different way.''
\item \textit{Agrees with} Section~\ref{sec:res:size-decides:parallel}, \emph{Where the tactical time goes: hand-built and suggested plans are close, with a small edge to Suggest in carrying the plan out}: The time a tactical suggestion frees is spent elsewhere while a plan runs. P36 describes that dead time from the commander's side: ``I've now got to wait for this thing to get into place'', and uses it to plan for failure.
\item \textit{Agrees with} Section~\ref{sec:limitations}, \emph{Limitations}: Others said the pace was unreal; P36 agrees (``in terms of the pace probably not so much'') and says what sets the real pace: ``if you're going at 3k an hour then that's going to take you a long time''.
\end{itemize}

\paragraph{For the parallel-work result, and the training discussion.}

\begin{quote}\small
``If I've sent that off it allows me to scan over here and then go back and check the answer over there. [...] That's what's going to be required if you're doing a drone swarm the whole point of it is to make sure that you have got [...] the key thing for me in all of this is not necessarily trusting the AI it is the skills that a human is going to need to be able to exploit all of this [...] that's what air traffic controllers do and they're actually testing it for it and their pay reflects that because it's not a particularly common skill.''\par\hfill--- P36 (think-aloud, session 2)
\end{quote}

\paragraph{Read against the results.}\phantomsection\label{sec:qual:practitioner:against-results}

\textcolor{blue}{[Claude: The old gathered list of links. The synthesis map at the start of the results now does this job; delete the list if you agree.]}


\paragraph*{Agrees with the results (24).}
\begin{itemize}[leftmargin=1.2em,itemsep=1pt]\small
\item Section~\ref{sec:res:size-decides:parallel}, from the point in Section~\ref{sec:qual:practitioner:pt-whole-board}: Their aim is throughput across the board (``get as many across as quickly as I can to complete all of the tasks''), the same reason the cohort gives for starting Suggest and moving on.
\item Section~\ref{sec:qual:hierarchy:cross-mission}, from the point in Section~\ref{sec:qual:practitioner:pt-whole-board}: Operators overrode the cards to act on what one mission's card cannot see. P36 says why in one line: ``success is generated across all of the tasks rather than just an individual task''.
\item Section~\ref{sec:res:manual-cards:dots}, from the point in Section~\ref{sec:qual:practitioner:pt-whole-board}: Playing for every task shows in the cohort's hand builds: on 44\% of decisions they sit below both cards, mostly by leaving off the card's spare, which keeps drones for the next mission.
\item Section~\ref{sec:res:size-decides:speed}, from the point in Section~\ref{sec:qual:practitioner:pt-quicker-by-hand}: The card saves time only on big three-type jobs. Three minutes into their first session P36 finds it out: ``Probably quicker just to hit three'', then ``I probably should have done that one manually. It would have been quicker.''
\item Section~\ref{sec:res:size-decides:both}, from the point in Section~\ref{sec:qual:practitioner:pt-branching-paths}: Small jobs by hand, big jobs to the assistant (cohort: card 56\% to 75\%, Suggest 54\% to 91\%). P36 gives both halves of the reason: on a simple job ``i know the answer based on what i'm seeing visually'', and a failure needs no re-plan (``you don't really need to change the story''); on a complex one ``my option scope is broader'', so ``let them start on that task and then intervene when i need to''. Their own log follows it (card 33\% to 70\%, Suggest 79\% to 100\%).
\item Section~\ref{sec:qual:trust-unchecked}, from the point in Section~\ref{sec:qual:practitioner:pt-branching-paths}: ``Give it a quick once over and go that looks about right ... whether or not it is right of course it's another question entirely'': the brief check the timing data show, with its limit named.
\item Section~\ref{sec:qual:hierarchy:one-plan}, from the point in Section~\ref{sec:qual:practitioner:pt-branching-paths}: P21 chained drones through each other's tasks for ``redundancy for free''. P36 does the same on simple missions (``every single drone is going to every single task so that if one of my drones fails ... my other drones are picking up the task''), which is why those plans need no checking afterwards.
\item Section~\ref{sec:res:workload:raw-vs-adjusted}, from the point in Section~\ref{sec:qual:practitioner:pt-resource-rich}: A busier moment does not push people onto the assistants. P36 requests a tactical suggestion at a quiet moment, ``because there's nothing else that's coming in''.
\item Section~\ref{sec:qual:new-jobs:queue-order}, from the point in Section~\ref{sec:qual:practitioner:pt-reserve}: Operators open the oldest waiting mission well above chance. P36: ``this one's been there for the longest i've got nothing in there that's going to change my prioritization''.
\item Section~\ref{sec:qual:new-jobs:queue-order}, from the point in Section~\ref{sec:qual:practitioner:pt-critical-first}: Operators lean towards the more critical waiting mission ($\times$1.45 per category, on top of oldest-first). P36 goes further and spends score on it: ``i've run out of blue so that's gonna hurt my score but ... i'd rather do the critical tasks if this is a real life scenario''.
\item Section~\ref{sec:qual:hierarchy:cascade}, from the point in Section~\ref{sec:qual:practitioner:pt-strategic-trust}: P19 put the only remedy for failures with the strategic tier. P36 reaches the same point in mid-sentence, and it reverses their reason for trusting the cards: ``they're probably not actually are they if you fuck up the allocation at the start of it then you're screwed going into the tactical''.
\item Section~\ref{sec:res:size-decides:parallel}, from the point in Section~\ref{sec:qual:practitioner:pt-physics-pace}: The time a tactical suggestion frees is spent elsewhere while a plan runs. P36 describes that dead time from the commander's side: ``I've now got to wait for this thing to get into place'', and uses it to plan for failure.
\item Section~\ref{sec:limitations}, from the point in Section~\ref{sec:qual:practitioner:pt-physics-pace}: Others said the pace was unreal; P36 agrees (``in terms of the pace probably not so much'') and says what sets the real pace: ``if you're going at 3k an hour then that's going to take you a long time''.
\item Section~\ref{sec:qual:trust-unchecked}, from the point in Section~\ref{sec:qual:practitioner:pt-staff-process}: Plans are accepted after a brief look. P36 explains why that can be rational: answers from staff are taken ``as long as you have trust in how they've come up with those answers'', not re-derived.
\item Section~\ref{sec:qual:trust-unchecked:cant-check}, from the point in Section~\ref{sec:qual:practitioner:pt-staff-process}: Several operators wanted to know how the assistants decide. For P36 that is the basis of trust: a commander can take the staff's answers because ``you'd have worked up in those levels so you kind of understand how those answers are generated''. No operator has done the assistant's job.
\item Section~\ref{sec:res:size-decides:speed}, from the point in Section~\ref{sec:qual:practitioner:pt-staff-process}: The card pays where a job is too big to do quickly by hand. P36 gives the same reason for having staff at all: ``I couldn't work through the ground the enemy, the enemy's plan, the likelihood of my plan succeeding, the logistics behind it it would take me forever which is why your staff do it''.
\item Section~\ref{sec:qual:trust-unchecked:cant-check}, from the point in Section~\ref{sec:qual:practitioner:pt-trust-but-verify}: ``With AI it should always be open to challenge'': the same ask as the operators who wanted to know how the assistants decide.
\item Section~\ref{sec:res:size-decides:parallel}, from the point in Section~\ref{sec:qual:practitioner:pt-play-catch-up}: Handing over frees the operator to work elsewhere: ``If I've sent that off it allows me to scan over here and then go back and check the answer over there'', which is where the cohort's Suggest time goes.
\item Section~\ref{sec:qual:trust-unchecked}, from the point in Section~\ref{sec:qual:practitioner:pt-play-catch-up}: Brief checking, whatever the plan size. P36 names its price: ``if I just said 'go nuts' when something goes wrong I've got to play catch up'', because the plan is not in your head.
\item Section~\ref{sec:qual:failure:blind}, from the point in Section~\ref{sec:qual:practitioner:pt-failure-rate}: Operators facing random failures with no numbers felt they were deciding blind. P36 agrees the numbers should be shown (``great to see the information'') and adds that people then need training to act on them: ``there's always going to be a failure rate I'm just trying to make it as low as reasonable''.
\item Section~\ref{sec:res:size-decides:tactical}, from the point in Section~\ref{sec:qual:practitioner:pt-resource-cascade}: Handing over at one tier goes with handing over at the other (Suggest $\times$2.1 after a card, same person, same size). P36 explains the coupling: the top allocation sets the lower tier's risk (``I told you need more resource when I told you this is my plan ... you're carrying risk''). Their own card-then-Suggest gap was close to zero, so the link is the cohort's, not theirs.
\item Section~\ref{sec:res:manual-cards:dots}, from the point in Section~\ref{sec:qual:practitioner:pt-resource-cascade}: Over-allocating at the top: ``that's just wasted resource''. In their own builds the cohort mostly trimmed the card's spare drone.
\item Section~\ref{sec:qual:hierarchy:cross-mission}, from the point in Section~\ref{sec:qual:practitioner:pt-resource-cascade}: Operators overrode the cards to weigh other missions. P36 says it as a commander: ``I've got five tasks that I need to do so I'm balancing these five'', and ``I've got a critical one over here that I'm allocating more to''.
\item (``Assistants that check rather than decide'', not in the paper), from the point in Section~\ref{sec:qual:practitioner:pt-red-team}: P11 wanted an assistant that checks their own plan. P36's checker tests the other agents' data instead, so a bad input shows up before the answers built on it.
\end{itemize}
\paragraph*{Qualifies the results (10).}
\begin{itemize}[leftmargin=1.2em,itemsep=1pt]\small
\item Section~\ref{sec:res:size-decides:both}, from the point in Section~\ref{sec:qual:practitioner:pt-resource-rich}: The size rule bends with slack: this mission is ``relatively simple'', but ``I'm resource rich at the moment so we'll go nuts I'll start suggesting''.
\item Section~\ref{sec:res:manual-cards:dots}, from the point in Section~\ref{sec:qual:practitioner:pt-reserve}: Most of the cohort's hand builds sit at or below the cards (44\% below both). P36 adds a spare on purpose (``minimum two we're going to go three just never deploy without a reserve''), yet their builds land where the cohort's do: on average 0.4 of a drone under the Conservative card, against a median of 0.6 under.
\item Section~\ref{sec:res:tier-trust:gap}, from the point in Section~\ref{sec:qual:practitioner:pt-strategic-trust}: The cohort trusts the Tactical Assistant more; P36 is one of the exceptions, and their ratings match what they say (``I trust the strategic slightly more'': 5.9 against 5.8 on the 1--7 scale, where the cohort's median strategic rating is 4.5). They hesitated less over the cards only because ``all i was looking at there was have i got a reserve allocated''.
\item Section~\ref{sec:res:workload:raw-vs-adjusted}, from the point in Section~\ref{sec:qual:practitioner:pt-physics-pace}: A busier queue did not push anyone onto the assistants. P36 puts the real pressure elsewhere: ``physics remains the bounding factor'', and the load becomes waiting and asking ``what am I going to do if this thing fails''. ``Your mental load is still loaded it's just loaded in a slightly different way.''
\item Section~\ref{sec:res:workload:said-did}, from the point in Section~\ref{sec:qual:practitioner:pt-tempo-echelon}: Nine people said they lean on the assistant when busy, yet their reliance did not rise with the queue. P36 puts the busy in a real hierarchy: ``the higher up you are the further out'', where a new item is ``just like a routine routine task''; at the lower tier ``10 things in front of me as it's all come down the chain''. The game gave one person both tiers.
\item Section~\ref{sec:qual:trust-unchecked}, from the point in Section~\ref{sec:qual:practitioner:pt-trust-but-verify}: This is from someone who kept every consulted tactical plan unedited. In P36's account verifying happens around the plan, not in a long look at it: take it to the people who will carry it out (``go away and work out the plan'') and keep it ``open to challenge''. Beside that sits ``trust in the fact it's done the maths behind it and has allocated enough platforms'', which is about all the cohort's 2--4 seconds of checking can cover.
\item Section~\ref{sec:qual:failure:blind}, from the point in Section~\ref{sec:qual:practitioner:pt-failure-rate}: Several operators asked for failure numbers. P36 wants them with context (``there's got to be something where it shows a bit of context to it''), because a bare ``40\% failure chance'' reads as ``no chance'' when ``the failure rate of doing nothing is 100\%''.
\item Section~\ref{sec:res:personal-style:not-disposition}, from the point in Section~\ref{sec:qual:practitioner:pt-shared-risk}: Prior experience did not predict style here. P36 expects it to matter most when the stakes are real: their own ``looking at a glance yes it's allocated what I needed ... cool let's start running with it'' comes from experience, where a new platoon commander would be ``foolhardy'' or ``almost paralysed''. A game with no lives on the line hides that difference.
\item Section~\ref{sec:qual:hierarchy:one-plan}, from the point in Section~\ref{sec:qual:practitioner:pt-silos}: Several operators wanted one plan held in mind across the tiers. P36 sees narrow context as normal (``you as a g2'er are will only see your part of the plan'') and fine ``as long as in each of them it's gone through the analytical process'', and adds that this game's tasks are simple enough not to need more: a direct tension with those operators.
\item Section~\ref{sec:qual:hierarchy:one-plan}, from the point in Section~\ref{sec:qual:practitioner:pt-one-up-two-down}: The operators who wanted one shared plan are answered by a different design: each level gets ``the understanding they're required to do that task'', and ``it doesn't matter how it docks into everything else''.
\end{itemize}
\paragraph*{Goes beyond the results (8).}
\begin{itemize}[leftmargin=1.2em,itemsep=1pt]\small
\item Section~\ref{sec:qual:new-jobs:queue-order}, from the point in Section~\ref{sec:qual:practitioner:pt-tempo-echelon}: P27 and P32 asked for an assistant that says what to open next. P36's lower tier is where that job would sit: ``I've got to do this thing now then that thing then that thing then that thing''.
\item Section~\ref{sec:qual:who-decides:loop}, from the point in Section~\ref{sec:qual:practitioner:pt-trust-but-verify}: Some operators want a human in the loop even one who cannot check. P36's check is not one person looking harder but a plan socialised so that ``somebody is checking my thinking and has the option to challenge me''.
\item Section~\ref{sec:res:personal-style:who-vs-what}, from the point in Section~\ref{sec:qual:practitioner:pt-play-catch-up}: Who decides predicts handover far better than the mission does. P36 draws the operational conclusion: ``the key thing for me in all of this is not necessarily trusting the AI it is the skills that a human is going to need to be able to exploit all of this'', a niche skill set like air traffic control, paid accordingly.
\item (``Lower-reliance operators attempt fewer recoveries'', not in the paper), from the point in Section~\ref{sec:qual:practitioner:pt-failure-rate}: Lower-reliance operators attempted fewer recoveries. P36's rule cuts the other way: in search and rescue ``you're always gambling with someone's life'', so you act and ``chip away at it''.
\item Section~\ref{sec:limitations}, from the point in Section~\ref{sec:qual:practitioner:pt-shared-risk}: ``It's a game'' was a common caveat. P36 says what changes with a life at stake (``I imagine you're going to see me less decisive if someone's gonna die at the end''), and that a commander who no longer shares the risk (``I might die alongside them'') loses what used to temper the order.
\item Section~\ref{sec:qual:hierarchy:one-plan}, from the point in Section~\ref{sec:qual:practitioner:pt-silos}: Operators wondered how the two assistants fit together. P36's model is staff cells that each see their part and first ``come together'' at the planning brief.
\item Section~\ref{sec:qual:hierarchy:one-plan}, from the point in Section~\ref{sec:qual:practitioner:pt-one-up-two-down}: How should the tiers fit? From doctrine: ``you understand two up in the same way as you plan two down'', each task comes with the part of the level above that it needs, and an agent reaching outside its lane is told ``that's not your job''.
\item Section~\ref{sec:res:flaw-odd:visibility}, from the point in Section~\ref{sec:qual:practitioner:pt-red-team}: The tactical flaw (a wrong ``No workable plan'') went almost unseen: only 4 people reacted to one aloud. P36's answer does not rely on the operator: ``an agent whose only job is to go and red team [adversarially check] all of your outputs very quickly'', checking ``the data that they've based their answers on''.
\end{itemize}

\begin{options}
\optshead{Options for this finding}

\optsec{How typical are they?} {\footnotesize [pass 5: numbers only]}

\noindent Your note: We will say that our expert was a good performer with a slightly higher workrate than others, prempting cards more, but was otherwise fairly typical (insert a couple numbers into here)

\optsec{How they played: for the whole board, not the mission} {\footnotesize [pass 5: for the paper]}

\noindent Your note: We'll briefly quote the second part I think

\noindent Highlighted quotes not used above:\begin{quotelist}
\item ``so generally speaking all I'm trying to do here is get as many across as quickly as I can to complete all of the tasks managing my resources that I've got left here so I'm running a little bit low on blue over there and then just double checking on here so I've got the redundancy in there''\,---\,P36 (think-aloud, session 1)\par\usernote{Your comment: Minor}
\end{quotelist}

\optsec{How they played: simple by hand, complex to the assistant, and why} {\footnotesize [pass 5: for the paper]}

\noindent Your note: This one is a new idea that we can talk about in more depth in the paper: More complicated tasks can have a more branching failure space. This is not something we cover in this work, but could be an important factor in future works.

\optsec{How they played: the rule bends when the reserve is full} {\footnotesize [pass 5: extra: may make it in]}

\noindent Your note: minor, may make it in

\optsec{How they played: oldest first, and never without a reserve} {\footnotesize [pass 5: for the paper]}

\noindent Your note: Short note - they identified an classified a rule ``never without a reserve''. Links to the consistency factor in training.

\optsec{How they played: critical tasks over score} {\footnotesize [pass 5: for the paper]}

\noindent Your note: good

\optsec{The task: the pressure lands on the lower tier} {\footnotesize [pass 5: for the paper]}

\noindent Your note: maybe include
\end{options}

\clearpage
